\chapter{System Design: Preprocessing Pipeline and Encrypted Division Protocol}
\chaptermark{System Design}
\HRule\\[2cm]
\label{Chapter4}
\lhead{\emph{\textbf{Chapter 4:} System Design}}

\begin{spacing}{1.6}

\section{System Architecture}

\subsection{Privacy Model and Threat Model}

The system is designed under an \textit{honest-but-curious} adversary model:
the detection server follows the protocol specification correctly but may
attempt to infer information about individual data points from the ciphertexts
it receives, the intermediate computations it performs, and the values it
returns to the client.

The privacy guarantee enforced throughout all protocol phases is as follows:

\begin{itemize}
  \item The \textbf{data owner (client)} holds the raw plaintext data, the
  secret decryption key $\mathbf{sk}$, and all final anomaly classification
  decisions. The client is the only entity that can decrypt any ciphertext.

  \item The \textbf{detection server} holds the public encryption key
  $\mathbf{pk}$, the relinearization keys $\mathbf{rlk}$ (used to reduce
  degree-2 ciphertexts after multiplication back to degree-1), and the Galois
  rotation keys $\mathbf{galk}$ (used for slot vector rotations). The server
  holds encrypted ciphertexts and performs all HE computation but
  \textbf{never} holds $\mathbf{sk}$.

  \item No individual data point is ever revealed to the server at any protocol
  step. The server only computes on ciphertexts.

  \item During the inversion rounds, the client receives only aggregate model
  parameters: $N$ (the number of training rows) and $\sigma_j^2$ (per-feature
  variance). These are global statistics over the entire training model 
  functionally equivalent to the stored parameters of a deployed ML model 
  not individual data points. Their disclosure does not expose any individual
  training observation.
\end{itemize}

\subsection{Three-Process Architecture}

The complete system is implemented as three independent processes:

\textbf{Process 1 --- Detection Server (FastAPI):}
Receives encrypted ciphertexts from the client via HTTP REST endpoints. Performs
all HE computation: slot sums, mean computation, variance derivation, PCA
fitting (in plaintext on the component matrix $\mathbf{V}$), and anomaly score
computation. Maintains a \texttt{ServerState} object that persists encrypted
intermediate values (training ciphertexts, computed means, inverse variances,
test ciphertexts) across multiple HTTP requests. The server holds only
$\mathbf{pk}$, $\mathbf{rlk}$, and $\mathbf{galk}$ (never $\mathbf{sk}$).
Score computation runs asynchronously in a background thread to prevent HTTP
timeout for large test sets.

\textbf{Process 2 --- Data Owner Client (embedded in Dashboard):}
Encrypts data using $\mathbf{pk}$, participates in inversion rounds (decrypts
aggregate parameters using $\mathbf{sk}$, computes exact reciprocals in
plaintext, re-encrypts using $\mathbf{pk}$), receives encrypted score
ciphertexts from the server, decrypts them using $\mathbf{sk}$, applies the
anomaly threshold, and computes evaluation metrics. All decryption operations
occur exclusively on the client.

\textbf{Process 3 --- Monitoring Dashboard (Flask + SocketIO):}
A live webbased monitoring interface that embeds the client logic, eliminating
the need for a separate client terminal. Browser buttons communicate via
SocketIO events to the dashboard backend, which spawns background threads to
execute client protocol operations. The dashboard displays real time protocol
progress, per round latency, score distributions, and classification metrics.

\subsection{Complete Module Structure}

\end{spacing}
\begin{table}[!ht]
\centering
\caption{Complete System Modules and Their Responsibilities}
\label{tab:modules}
\begin{tabular}{p{5cm}p{8cm}}
\toprule
\textbf{Module} & \textbf{Responsibility} \\
\midrule
\texttt{key.py} & SEAL context initialization, key generation (pk, sk, rlk, galk), key persistence \\
\texttt{shared/config.py} & Centralized configuration: network ports, file paths, SEAL parameters, protocol constants \\
\texttt{shared/protocol.py} & Base64 serialization/deserialization of SEAL ciphertexts for HTTP JSON transport \\
\texttt{client/encrypt.py} & Column-major training ciphertext creation; row-major test ciphertext creation; scaler application \\
\texttt{client/decrypt.py} & Scalar single-value decryption; vector multi-value decryption \\
\texttt{client/invert.py} & Exact $1/x$ computation in plaintext and CKKS re-encryption \\
\texttt{client/decide.py} & Threshold classification; AUC, F1, Precision, Recall, confusion matrix computation \\
\texttt{client/http\_client.py} & HTTP client implementing all SAD and PCA protocol rounds over REST \\
\texttt{server/algo.py} & HE-SAD: slot sums, mean computation, sigma² computation, SAD score computation \\
\texttt{server/pca\_algo.py} & HE-PCA: slot sums, mean computation, PCA fitting, optimal $k^*$ selection, reconstruction error scoring \\
\texttt{server/http\_server.py} & FastAPI REST server, ServerState machine, async score computation, results management \\
\texttt{dashboard/dashboard\_app.py} & Flask + SocketIO dashboard; embedded client controller; browser button event handlers \\
\texttt{dashboard/templates/dashboard.html} & Live monitoring UI: buttons, progress bar, round status, latency charts, metrics display \\
\texttt{sad\_baseline.py} & Plaintext SAD reference implementation for performance comparison \\
\texttt{pca\_baseline.py} & Plaintext PCA reference implementation without StandardScaler \\
\bottomrule
\end{tabular}
\end{table}

\begin{spacing}{1.6}
    
\end{spacing}

\section{Dataset-Agnostic Preprocessing Pipeline}

\subsection{Design Philosophy}

The preprocessing pipeline converts raw multivariate data from any domain into
CKKS-compatible normalized form. Three constraints drive the design:

\textbf{CKKS compatibility:} All feature values must lie within $[-1, 1]$ after
normalization. CKKS encodes plaintext vectors as polynomial coefficients scaled
by $2^{40}$. Values outside $[-1, 1]$ cause coefficient overflow, and values
far from zero amplify noise relative to signal. The normalization step is
therefore not merely a conventional ML preprocessing step but a CKKS correctness
requirement.

\textbf{Information separation:} All normalization parameters must be computed
exclusively from normal training data. If min/max values were computed from
data including attack rows, extreme attack-induced values would expand the
normalization range, compressing all normal values into a narrow band near zero
and reducing detection sensitivity. The pipeline enforces strict separation
between training statistics and test data.

\textbf{Feature type awareness:} Datasets contain heterogeneous feature types
--- continuous analog sensor readings and discrete binary actuator states. These
require different normalization strategies to maintain appropriate value ranges
and to ensure each feature type contributes meaningfully to anomaly scoring.

\subsection{Stage 1: Structural Cleaning}

The first stage removes all metadata columns that carry no sensor signal and
would corrupt the numeric feature matrix:

\begin{itemize}
  \item \textbf{Timestamp columns} (\texttt{Timestamp} in SWaT,
  \texttt{DATETIME} in BATADAL): Date-time strings cannot be meaningfully
  encoded as real valued features and carry no anomaly discriminating information.
  \item \textbf{Label columns} (\texttt{Normal/Attack} in SWaT,
  \texttt{ATT\_FLAG} in BATADAL): Ground truth labels are retained separately
  for evaluation but excluded from all feature processing, ensuring the system
  remains unsupervised.
\end{itemize}

After cleaning, SWaT retains 51 feature columns and BATADAL retains 43,
consisting purely of numeric sensor and actuator readings.

\subsection{Stage 2: Missing Value Resolution}

The second stage fills gaps in the numeric feature matrix using a three-pass
imputation strategy:

\textbf{Pass 1 --- Forward fill:} The last valid observation for each feature
is propagated forward to fill subsequent gaps. This is the physically appropriate
assumption for time series sensor data: a missing reading most likely reflects
a sensor reporting its last known state, since physical quantities such as water
level, pressure, and flow rate change continuously and cannot jump discontinuously.

\textbf{Pass 2 --- Backward fill:} Any gaps remaining at the beginning of the
recording period (where no prior valid observation exists to propagate forward)
are resolved by propagating the next valid observation backward.

\textbf{Pass 3 --- Column mean fill:} As a last resort, any gaps persisting
after both passes are filled with the column mean. In practice, both SWaT and
BATADAL required only the first two passes, with post resolution verification
confirming zero remaining missing values.

A programmatic check (\texttt{np.isnan(X).sum() == 0}) confirms successful
imputation before proceeding.

\subsection{Stage 3: Feature Type Classification}

Each feature is automatically classified as continuous or binary based on its
observed value distribution in the training data:

\textbf{Continuous features} are real-valued analog sensor readings that vary
smoothly over time. Examples include:
\begin{itemize}
  \item Flow rates (FIT/F\_PU sensors): measure volumetric flow in litres per
  second, ranging from near-zero to several hundred.
  \item Water levels (LIT/L\_T sensors): measure absolute level in millimetres
  or centimetres.
  \item Pressures (PIT/P\_J sensors): measure differential or absolute pressure
  in appropriate units.
\end{itemize}

\textbf{Binary features} are discrete indicators recording the on/off or
open/closed state of digital actuators. Examples include:
\begin{itemize}
  \item Pump on/off signals (P/S\_PU sensors): 0 = off, 1 = on.
  \item Motorized valve positions (MV/S\_V sensors): 0 = closed, 1 = open.
\end{itemize}

SWaT contains 26 continuous and 25 binary features; BATADAL contains 28
continuous and 15 binary features.

\subsection{Stage 4: CKKS-Compatible Normalization}

\subsubsection{Continuous Feature Normalization}

Each continuous feature $x_i$ is linearly scaled to $[-1, 1]$ using the
minimum and maximum values observed across all \textit{normal} training rows:
\begin{equation}
  x'_i = \frac{2(x_i - x_{\min})}{x_{\max} - x_{\min}} - 1
  \label{eq:normalization}
\end{equation}
where $x_{\min}$ and $x_{\max}$ are computed from normal training rows only.
The scaled value equals $-1$ at the minimum observed normal value, $+1$ at the
maximum, and $0$ at the midpoint of the normal operating range.

In cases where $x_{\max} = x_{\min}$ (a constant feature), the denominator is
replaced by $1$ to avoid division by zero, and the feature is assigned value $0$.
The per-feature $(x_{\min}, x_{\max})$ pairs are serialized to
\texttt{scaler.npy} for consistent application to test data.

\subsubsection{Binary Feature Normalization}

Binary actuator state features are remapped from their original $\{0, 1\}$
representation to $\{-1, +1\}$:
\begin{equation}
  x'_i = 2x_i - 1
  \label{eq:binary_norm}
\end{equation}

Two considerations motivate this remapping:

\textit{Uniform feature space:} Remapping brings binary features into the same
$[-1, +1]$ value space as normalized continuous features. When all $d$ features
are packed into a single CKKS ciphertext slot vector, each feature occupies one
slot. If binary features retained their $\{0, 1\}$ range while continuous
features span $[-1, 1]$, the binary features would have systematically half the
magnitude, distorting the anomaly score contributions.

\textit{Symmetric anomaly contribution:} In the SAD score formula, the squared
deviation of a binary feature from its mean $\mu_j \in (-1, 1)$ is computed.
With $\{-1, +1\}$ values, both the off-state ($-1$) and the on-state ($+1$)
contribute nonzero squared deviations from any mean, ensuring that actuator
state anomalies (an unexpected pump activation or valve closure) register in
the anomaly score.

\subsection{Stage 5: Output Verification}

Five automated checks guard against silent preprocessing failures:
\begin{enumerate}
  \item Zero missing values: \texttt{np.isnan(X).sum() == 0}
  \item Continuous range: all values within $[-1.01, 1.01]$ (1\% tolerance for
  floating-point rounding at boundary values)
  \item Binary value set: each binary column contains only $\{-1, +1\}$
  \item Row count preservation: output row count equals input row count
  \item Column count: exactly 51 (SWaT) or 43 (BATADAL) columns
\end{enumerate}

Both datasets passed all five checks, producing verified output files used as
direct input to the CKKS encryption pipeline.

\begin{table}[H]
\centering
\caption{Preprocessing Pipeline Summary}
\label{tab:preproc}
\begin{tabular}{p{0.5cm}p{3.5cm}p{4.5cm}p{4.5cm}}
\toprule
\textbf{\#} & \textbf{Stage} & \textbf{SWaT} & \textbf{BATADAL} \\
\midrule
1 & Structural Cleaning & Remove Timestamp, Normal/Attack & Remove DATETIME, ATT\_FLAG \\
2 & Missing Values & Fwd $\to$ Bwd $\to$ Mean fill & Fwd $\to$ Bwd $\to$ Mean fill \\
3 & Feature Classification & 26 continuous, 25 binary & 28 continuous, 15 binary \\
4a & Continuous Norm. & Min-max to $[-1,1]$ from normal rows & Min-max to $[-1,1]$ from normal rows \\
4b & Binary Norm. & $\{0,1\} \to \{-1,+1\}$ & $\{0,1\} \to \{-1,+1\}$ \\
5 & Verification & 5 checks, all passed & 5 checks, all passed \\
\bottomrule
\end{tabular}
\end{table}

\section{CKKS Encoding Strategies}

Two distinct encoding strategies are used for training data and test data,
each optimised for the operations performed on them.

\subsection{Column-Major Training Encoding}

Training data is encoded one ciphertext per feature, packing all $N$ training
row values for that feature into the SIMD slots:
\begin{equation}
  \mathrm{ct}_j^{\text{train}} = \mathrm{Enc}_\mathbf{pk}
  ([x_{0j}, x_{1j}, x_{2j}, \ldots, x_{N-1,j}])
  \quad \text{for } j = 1, \ldots, d
  \label{eq:colmajor}
\end{equation}

This encoding enables the computation of the feature column sum using a
binary rotation tree: summing $N$ slot values requires only $\lceil\log_2 N
\rceil$ rotate-and-add steps. For $N = 5000$ training rows, this requires 13
rotations rather than 4999 sequential additions, reducing noise accumulation
and computation time by orders of magnitude.

The column sum $\sum_{i=0}^{N-1} x_{ij}$ (present in all slots after the
rotation tree) is used to compute the feature mean $\mu_j = \frac{1}{N}
\sum_i x_{ij}$ after division by $N$ via the multi-round protocol.

\subsection{Row-Major Test Encoding}

Test data is encoded one ciphertext per test observation, packing all $d$
feature values into slots $0$ through $d-1$, with slots $d$ through $n/2-1$
zero-padded:
\begin{equation}
  \mathrm{ct}_i^{\text{test}} = \mathrm{Enc}_\mathbf{pk}
  ([x_{i0}, x_{i1}, \ldots, x_{i,d-1}, 0, 0, \ldots, 0])
  \quad \text{for } i = 1, \ldots, N_\text{test}
  \label{eq:rowmajor}
\end{equation}

This encoding enables anomaly score computation for observation $i$ in a single
HE evaluation: subtract the mean vector (sub\_plain), square element-wise
(square), multiply by the inverse variance vector (mult\_plain), and sum all
$d$ slot values using $\lceil\log_2 d\rceil$ rotate-and-add steps. For
$d = 51$ features, only 6 rotations are required.

The zero padding in slots $d$ through $n/2-1$ contributes zero to the slot
sum, so the anomaly score for the row is correctly accumulated in slot 0 of
the result ciphertext.

\section{The Encrypted Division Problem}

\subsection{Problem Statement}

The CKKS scheme natively supports ciphertext-ciphertext addition and
multiplication, and ciphertext-plaintext addition and multiplication. It does
not natively support division. This limitation is fundamental: division is
a non-polynomial function and cannot be computed exactly using a finite number
of multiplications.

Both SAD and PCA require division at critical protocol steps:
\begin{itemize}
  \item \textbf{Computing the feature mean:} $\mu_j = \frac{1}{N} \sum_i
  x_{ij}$ requires dividing the column sum by $N$.
  \item \textbf{SAD variance normalization:} $s(\mathbf{x}) = \sum_j
  \frac{(x_j-\mu_j)^2}{\sigma_j^2}$ requires dividing by $\sigma_j^2$ for each
  of $d = 51$ features.
\end{itemize}

Without division, these fundamental statistics cannot be computed, and the
anomaly detection algorithms cannot be faithfully implemented under encryption.

Three approaches to encrypted division were evaluated: Newton-Raphson iterative
inversion, Chebyshev polynomial approximation, and a novel multi-round
interactive protocol.

\subsection{Approach 1: Newton-Raphson Iterative Inversion}

The Newton-Raphson method for computing $1/a$ is the iterative formula:
\begin{equation}
  y_{n+1} = y_n(2 - a \cdot y_n)
  \label{eq:nr}
\end{equation}

Under the convergence condition $|1 - a \cdot y_0| < 1$, this sequence
converges quadratically to $1/a$ starting from the initial guess $y_0$. Each
iteration requires two multiplications (computing $a \cdot y_n$ and
$y_n \cdot (2 - a \cdot y_n)$), consuming 2 HE levels per iteration.

Applied to computing $1/\sigma_j^2$ for all 51 features simultaneously under
CKKS, the 51 variance values are packed into the 51 slots of a single
ciphertext, and the iteration runs element wise across all slots simultaneously.

\textbf{Experimental results on SWaT variance values:}

The SWaT training data produces feature variances $\sigma_j^2 \in [0.012,
0.45]$, spanning a $37\times$ range from smallest to largest. A single initial
guess $y_0$ must work for all 51 slots simultaneously:

\begin{table}[H]
\centering
\caption{Newton-Raphson Convergence on SWaT $\sigma^2 \in [0.012, 0.45]$}
\label{tab:nr}
\begin{tabular}{cccc}
\toprule
\textbf{Iter.} & \textbf{$y_0 = 4.33$ (1/median)} & \textbf{Mean Error} & \textbf{Status} \\
\midrule
0 & Initial guess & 59.8\% & Single guess for all 51 slots \\
1 & $y_1 = y_0(2-ay_0)$ & 187.2\% & Diverging \\
2 & $y_2 = y_1(2-ay_1)$ & 941.1\% & Fully diverged \\
3 & $y_3 = y_2(2-ay_2)$ & 5594.3\% & Unusable \\
\bottomrule
\end{tabular}
\end{table}

\textbf{Root cause of divergence:} The convergence condition requires
$|1 - a \cdot y_0| < 1$, equivalently $y_0 \in (0, 2/a)$. For a single
initial guess $y_0 = 4.33$ (the reciprocal of the median variance):
\begin{itemize}
  \item For the maximum variance feature ($\sigma^2 = 0.45$): $a \cdot y_0
  = 0.45 \times 4.33 = 1.95$, giving $|1 - 1.95| = 0.95 < 1$. Convergent
  but slow.
  \item For the minimum variance feature ($\sigma^2 = 0.012$): $a \cdot y_0
  = 0.012 \times 4.33 = 0.052$, giving $|1 - 0.052| = 0.948 < 1$.
  Technically convergent, but the $37\times$ range means many features
  are near the boundary of convergence while others diverge under the
  combined SIMD iteration.
\end{itemize}

Per slot initial guesses assigning $y_0^{(j)} = 1/\sigma_j^2$ to each slot
independently  would solve this problem, but CKKS does not support
slot-selective initial assignments in a way that avoids the underlying
arithmetic instability when all slots evolve under the same iteration formula.

\textbf{Conclusion:} Newton-Raphson inversion is not viable for encrypted
division across heterogeneous variance values.

\subsection{Approach 2: Chebyshev Polynomial Approximation}

A degree-$n$ Chebyshev polynomial is fitted offline to approximate $1/x$ on
the interval $[x_{\min}, x_{\max}]$ using Chebyshev nodes to minimize the
Runge phenomenon. The polynomial is then evaluated inside the HE computation
by substituting the ciphertext for $x$.

\begin{table}[H]
\centering
\caption{Chebyshev Approximation of $1/x$ on $[0.012, 0.45]$}
\label{tab:cheb}
\begin{tabular}{cccl}
\toprule
\textbf{Degree} & \textbf{Plaintext Error} & \textbf{HE Levels} & \textbf{HE Status} \\
\midrule
3 & 33.3\% & 2--3 & Executed; error unacceptable \\
5 & 19.1\% & 3--4 & Scale overflow \\
7 & 8.5\%  & 3--4 & Scale overflow \\
9 & 5.4\%  & 4--5 & Scale overflow \\
\bottomrule
\end{tabular}
\end{table}

\textbf{Analysis of results:}

At degree 3, the polynomial executes but produces a mean approximation error
of 33.3\% across the variance range. Substituting $1/\sigma_j^2$ with a value
33\% in error means the anomaly score for each feature deviates by 33\% from
the correct normalised squared deviation. This completely destroys the score
separation between normal and anomalous observations, making degree-3 Chebyshev
unusable.

At degrees 5, 7, and 9, scale overflow occurs. The root cause: each Chebyshev
basis polynomial $T_k(x)$ must be evaluated by multiplying the plaintext
polynomial coefficients with the current ciphertext level. The plaintext was
encoded at the global scale $2^{40}$, but after several rescaling operations
the ciphertext operates at a smaller effective scale determined by the remaining
prime size. Multiplying a plaintext encoded at $2^{40}$ by a ciphertext at a
reduced scale produces a product with scale far exceeding the current prime
capacity, causing overflow.

Encoding the plaintext at the ciphertext's current scale (\texttt{ct.scale()}
rather than the global \texttt{SCALE}) resolves the overflow but reduces the
polynomial evaluation precision, and degree 3 remains the only depth-feasible
option, with unacceptable 33.3\% error.

\textbf{Conclusion:} Chebyshev polynomial approximation is not viable for
encrypted division in this application.

\subsection{Approach 3: Multi-Round Interactive Protocol (Adopted)}

The multi-round interactive protocol resolves the encrypted division problem
by routing the computation of the reciprocal through the key holding data owner,
who can perform exact division in plaintext.

The protocol for computing $\mathrm{ct}_{1/a}$ from $\mathrm{ct}_a$ is:

\begin{enumerate}
  \item The server computes and sends $\mathrm{ct}_a$ to the client.
  \item The client decrypts using $\mathbf{sk}$: $a = \mathrm{Dec}_\mathbf{sk}
  (\mathrm{ct}_a)$.
  \item The client computes the exact reciprocal in plaintext:
  $r = 1.0 / a$ (IEEE 754 double precision).
  \item The client encodes $r$ and re-encrypts using $\mathbf{pk}$:
  $\mathrm{ct}_r = \mathrm{Enc}_\mathbf{pk}(r)$.
  \item The client sends $\mathrm{ct}_r$ to the server, which continues the
  protocol using $\mathrm{ct}_r$ in subsequent multiplications.
\end{enumerate}

\begin{table}[H]
\centering
\caption{Comparison of Division Approaches}
\label{tab:div}
\begin{tabular}{p{3cm}cccp{3.5cm}}
\toprule
\textbf{Method} & \textbf{Error} & \textbf{HE Levels} & \textbf{Works} & \textbf{Limitation} \\
\midrule
Newton-Raphson & Diverges & 6 per attempt & No & Wide variance range causes divergence \\
Chebyshev d=3 & 33.3\% & 2--3 & No & Error too large for anomaly detection \\
Chebyshev d$\geq$5 & $<$20\% & 3--5 & No & Scale overflow in HE \\
\textbf{Multi-Round} & \textbf{0\%} & \textbf{0} & \textbf{Yes} & None \\
\bottomrule
\end{tabular}
\end{table}

\textbf{Privacy analysis:} The protocol reveals to the client only the value
$a$ that was decrypted. In the SAD protocol, this is either:
\begin{itemize}
  \item $N = 5000$ (the number of training rows): reveals only the dataset
  size, which is a public protocol parameter, not a data point.
  \item $\sigma_j^2$ for $j = 1, \ldots, 51$: the feature variances over the
  training set. These are global statistics computed over all 5000 training
  rows, equivalent in sensitivity to the stored parameters of a trained
  scikit-learn model. They do not reveal any individual training observation
  or its label.
\end{itemize}

The protocol introduces two additional communication rounds but adds zero HE
multiplication levels, preserving the full 5-level budget for the actual anomaly
scoring computation.

\textbf{Conclusion:} The multi-round interactive protocol achieves exact
encrypted division with zero approximation error and zero HE depth cost, making
it the only viable approach for this application.

\end{spacing}
